\documentclass[a4paper,12pt]{article}
\usepackage[utf8]{inputenc}
\usepackage{geometry}
\usepackage{titlesec}
\usepackage{enumitem}
\usepackage{hyperref}
\usepackage{graphicx}
\usepackage{xcolor}

% Page margins
\geometry{left=25mm, right=25mm, top=25mm, bottom=25mm}

% Custom section formatting
\titleformat{\section}{\large\bfseries\color{blue}}{}{0em}{}[\titlerule]
\titleformat{\subsection}{\normalsize\bfseries}{}{0em}{}

\title{\textbf{Idea Submission: AI Unlocked}}
\author{Team Name: Devil2b \\ Lead Developer: Shashwat Pratap Singh}
\date{}

\begin{document}

\maketitle

\section{General Information}
\textbf{Institute Name:} IIT BHU, Varanasi \\
\textbf{Challenge Area:} Track 6: The Next Billion Builders \\
\textbf{Team Lead Email:} shashwatp.singh.cse24@itbhu.ac.in

\section{Project Title}
\textbf{Title:} StackStore: The AI-Native Runtime Manager \\
\textbf{Sub-title:} Instant, isolated development environments generated by AI -- no configuration required.

\section{The Team}
\textit{Tell us the story about your team. How did you decide to work together?} \\
\\
I am a student currently working on a custom OS-level virtualization engine for my academic research. While building the core isolation logic, I realized that the biggest bottleneck wasn't the kernel code—it was the complexity of configuring it. I decided to participate in AI Unlocked to build the "brain" for my engine: an AI layer that automates the complex setup . This project bridges my deep-tech systems idea with Microsoft's cutting-edge AI, aiming to solve "dependency hell" for good.

\section{The Concept}
\textit{What problem are you trying to solve, and how are you solving it?} \\

\textbf{The Problem:} Modern software development requires complex environment setups (Dockerfiles, dependencies, version managers). For students and new developers, this is a massive barrier to entry. Existing tools like GitHub Codespaces are expensive cloud solutions, while raw Docker is often too complex for beginners.

\textbf{The Solution:} StackStore is a local desktop application that uses AI to automate environment orchestration. Instead of requiring users to write configuration files manually, StackStore allows a user to simply input a GitHub repository URL. The integrated AI agent analyzes the source code, detects the language/framework (e.g., Python/Flask, Node.js/React), and automatically generates a secure, isolated runtime specification. The application then spins up this environment instantly on the user's local machine using the custom virtualization engine I am developing.

\section{Target Audience}
\begin{itemize}
    \item \textbf{Primary:} Computer Science students and hackathon participants who need to switch between tech stacks quickly without "polluting" their host OS.
    \item \textbf{Secondary:} Open-source maintainers who need to review Pull Requests (PRs) safely without setting up local environments manually.
    \item \textbf{Market Estimation:} With over 100 million developers on GitHub, a significant portion are students or junior developers who struggle with DevOps configuration.
\end{itemize}

\section{Personas}
\textbf{Persona 1: Riya (The Student Developer)} \\
Riya wants to try a new open-source project for her assignment. She usually spends 4 hours debugging "Python version mismatch" errors. With StackStore, she pastes the repo URL, and the AI sets up the correct version and dependencies in a sandbox instantly.

\textbf{Persona 2: Arjun (The Maintainer)} \\
Arjun maintains a popular library. He receives 10 PRs a week. Instead of manually configuring his laptop to test each one, he uses StackStore to spin up temporary, isolated environments for every PR, ensuring safety and speed.

\section{How It Works}
StackStore bridges the gap between raw source code and execution using a two-layer architecture facilitated by Microsoft Azure:

\begin{enumerate}
    \item \textbf{Input Analysis (The AI Layer):} The user provides a repository URL. The application sends key configuration files (e.g., \texttt{package.json}, \texttt{requirements.txt}) to \textbf{Azure OpenAI (GPT-4o)}.
    \item \textbf{Spec Generation:} The AI agent parses these files to understand the dependency tree. It generates a standardized "StackSpec"---a JSON file defining the necessary OS-level isolation (filesystem roots, network ports, env vars).
    \item \textbf{Execution (The OS Layer):} The StackStore engine reads this spec and uses OS-level virtualization (Linux Namespaces/cgroups) to create a pristine, sandboxed environment locally.
    \item \textbf{Self-Healing:} If the build fails, error logs are fed back to the AI, which suggests fixes to the configuration automatically.
\end{enumerate}

\section{Core Technologies}
\begin{itemize}
    \item \textbf{Microsoft Azure OpenAI Service:} Powers the code analysis and configuration generation.
    \item \textbf{StackStore Engine (C++/Rust):} Custom OS-level virtualization logic for lightweight isolation (Core innovation).
    \item \textbf{Electron/React:} Cross-platform desktop GUI for the user.
    \item \textbf{Azure Container Apps (Future):} For offloading heavy compilation tasks to the cloud.
\end{itemize}

\section{The Business Plan}
\begin{itemize}
    \item \textbf{Freemium Model:} "StackStore Personal" is free for students and open-source usage (local execution).
    \item \textbf{Pro Tier:} "StackStore Cloud" allows teams to share "live environment links" (stateful URLs) for collaborative debugging.
    \item \textbf{Enterprise:} On-premise deployment for companies needing secure, isolated sandboxes for proprietary code analysis.
\end{itemize}

\section{Additional Information}
\textit{Links or info you want to share:}
\begin{itemize}
    \item \textbf{GitHub Repo:} \href{https://github.com/PratapShashwat/DevRun}{PratapShashwat/DevRun}  (Project not started yet.) 
    \item \textbf{Future Roadmap:} I plan to implement "Natural Language Environment Editing," allowing users to say "Add a Redis database to this project" and have the AI update the configuration instantly.
\end{itemize}

\end{document}